\documentclass[11pt]{article}

\usepackage[utf8]{inputenc}
\usepackage[spanish]{babel}
\usepackage[T1]{fontenc}
\usepackage{amsmath, amssymb}
\usepackage{geometry}
\geometry{margin=2.5cm}
\usepackage{enumitem}
\usepackage{csquotes}
\usepackage{hyperref}
\usepackage[authoryear]{natbib}

\title{Portafolios de crecimiento robusto con costos de transacción:\\
una aproximación Wasserstein--Kelly informada por estados del mercado estadounidense}
\author{Juan Tomás Macías Yepes}
\date{\today}

\begin{document}
\maketitle{}

\section{Introducción}

\subsection{Contexto y antecedentes del problema}

%%% [C1] Dir. p.1 (Rockafellar y Uryasev / repetición de Kelly): YA APLICADO por el autor («por ejemplo, ... para este último» y Kelly (1956) solo una vez). Sin cambios.
Construir un portafolio consiste en decidir cómo repartir una cantidad de riqueza entre un
conjunto de activos riesgosos y un activo de bajo riesgo, antes de conocer los retornos que
esos activos producirán. Los enfoques más conocidos para esta decisión en sus formulaciones estáticas estándar -- el modelo media–varianza de \citet{markowitz1952} y los enfoques de optimización basados en VaR o CVaR  (por ejemplo,  \citet{rockafellaruryasev2000} para este último) -- priorizan compromisos entre retorno y riesgo o el control de pérdidas adversas. El criterio introducido por \citet{kelly1956}, en cambio, tiene como objetivo directo el crecimiento multiplicativo de la riqueza a largo plazo.
Formalizado por \citet{breiman1961}, responde
directamente a esa pregunta: maximiza la esperanza del logaritmo del factor de crecimiento
del portafolio, lo cual --bajo retornos independientes e idénticamente distribuidos-- maximiza
la tasa asintótica de crecimiento de la riqueza entre los portafolios de ponderaciones
constantes. El problema práctico es que la distribución de los retornos futuros es
desconocida; el portafolio de Kelly empírico se calcula sobre una distribución estimada a
partir de una ventana histórica finita, y el criterio logarítmico es conocido por ser sensible a
errores en esa estimación: un portafolio que parece óptimo sobre los escenarios pasados
puede estar sobreajustado a ellos y presentar un desempeño deficiente fuera de muestra \citep{li2023}.

La optimización robusta distribucionalmente (DRO) ofrece una respuesta a ese problema:
en lugar de optimizar contra la distribución empírica, se optimiza contra el peor caso dentro
de un conjunto de distribuciones plausibles construido alrededor de ella. Cuando ese
conjunto se define mediante una bola de Wasserstein, la protección incorpora además una
noción de geometría entre escenarios --trasladar probabilidad hacia un escenario similar es
más barato que trasladarla hacia uno muy distinto--, lo cual resulta natural cuando se quiere
comparar meses históricos usando tanto sus retornos como el contexto de mercado en que
ocurrieron. La combinación de Kelly con robustez distribucional de Wasserstein ya ha sido estudiada en la literatura: \citet{sunboyd2018} formularon el Kelly robusto distribucionalmente y
estudiaron sus propiedades de largo plazo; \citet{li2023} propuso un portafolio Wasserstein-Kelly
centrado en la distribución empírica y mostró mejoras de estabilidad fuera de muestra frente
al Kelly nominal; \citet{hsiehyu2024} incorporaron costos de transacción convexos y
condiciones de supervivencia al portafolio logarítmico robusto. Por separado, \citet{liu2022}
mostraron que la geometría utilizada para comparar los datos --no solo el tamaño del
conjunto de ambigüedad-- puede afectar de manera sustancial el comportamiento de un
modelo DRO.

\subsection{Vacío que se atiende}\label{sec:vacio}

Hasta donde alcanza la revisión bibliográfica realizada, no se ha identificado una comparación controlada que, dentro de un portafolio Wasserstein–Kelly con costos de transacción, contraste una geometría definida únicamente sobre retornos con una geometría que incorpore explícitamente variables de estado del mercado como el retorno agregado y el VIX. Se incorporan el VIX y el retorno del mercado como variables exógenas debido a la información que proporcionan sobre las condiciones actuales y las expectativas de los participantes del mercado. En particular, el VIX constituye una medida de la volatilidad esperada, construida a partir de los precios de opciones sobre índices, por lo que refleja las expectativas del mercado acerca de la volatilidad futura \citep{whaley2000}. Por otra parte, existe evidencia de una relación sistemática entre los retornos del mercado y la volatilidad, encontrándose una relación negativa entre los retornos inesperados y los 
%%% [C3a] Dir. p.2 (French et al.): se precisa que French et al. (1987) estudian volatilidad REALIZADA, distinta del VIX (volatilidad implícita prospectiva).
%%% Texto anterior: cambios inesperados en la volatilidad \citep{french1987}.
cambios inesperados en la volatilidad realizada \citep{french1987}.
%%% [fin C3a]
 
%%% [C3b] Dir. p.2 (French et al.): se presenta la elección como motivación a probar, no como evidencia de mejora ni de régimen predictivo.
%%% Texto anterior: En conjunto, estas variables permiten incorporar tanto información sobre la volatilidad esperada como sobre la dirección y condiciones recientes del mercado.
Esta evidencia no demuestra, por sí sola, que añadir el retorno de mercado y el VIX a la distancia de transporte mejore la optimización ni que exista un régimen predictivo; además, el VIX es una medida prospectiva de volatilidad implícita, mientras que \citet{french1987} estudian la volatilidad realizada. Por ello, la incorporación de estas variables se presenta como una motivación económicamente razonable que este estudio pone a prueba empíricamente, y no como un resultado supuesto.
%%% [fin C3b]
 

%%% [C2] Dir. p.2 (Sun y Zou): referencia ya incorporada por el autor; se corrige \citep -> \citet porque funciona como sujeto de la oración.
%%% Texto anterior: El estudio de \citep{sunzou2025} plantea
El estudio de \citet{sunzou2025} plantea
%%% [fin C2]
 una nueva clase de modelos coherentes de optimización de portafolios de Kelly basados en datos y en la métrica de Wasserstein, donde compara su desempeño con el modelo nominal de Kelly. Sin embargo, este no necesariamente reproduce distintos tipos de geometrías de transporte concretas, ni incorpora costos de transacción.  

%%% [C4] Dir. p.2 (redacción del vacío): ya reformulada por el autor; solo se corrige la concordancia («consisten» -> «consiste»).
%%% Texto anterior: El vacío identificado consisten en
El vacío identificado consiste en
%%% [fin C4]
 la falta de evidencia sobre si definir la distancia entre escenarios históricos
únicamente en los retornos de los activos, frente a anclarla también en el estado del mercado en que esos retornos ocurrieron, produce un compromiso distinto entre crecimiento, pérdidas, concentración y costos de transacción. Ninguna de las referencias anteriores compara directamente estas dos
geometrías bajo el mismo universo de activos, el mismo periodo, las mismas restricciones y
el mismo protocolo de backtest.

\subsection{Pregunta de investigación}

\begin{displayquote}
\textbf{¿Cómo afecta la incorporación de variables de estado del mercado en la geometría de transporte de un portafolio Wasserstein–Kelly con costos de transacción al desempeño del portafolio, en comparación con una geometría basada únicamente en los retornos de los activos?}
\end{displayquote}

\subsection{Objetivo general}

Evaluar el efecto de incorporar el retorno de mercado y el VIX en la geometría de transporte de un portafolio Wasserstein–Kelly con costos de transacción proporcionales sobre su desempeño fuera de muestra, en comparación con una geometría basada únicamente en los retornos de los activos, utilizando el criterio de Kelly nominal como referencia adicional y considerando crecimiento, pérdidas, concentración y costos de transacción.

\subsection{Objetivos específicos}

\begin{enumerate}
    \item \textbf{Establecer} el desempeño de referencia del modelo de Kelly nominal con costos de transacción proporcionales, bajo las restricciones de balance de cartera, límites de posición y supervivencia consideradas en el estudio.
\item \textbf{Formular} y validar el modelo Wasserstein–Kelly bajo dos geometrías de transporte: una definida únicamente sobre los retornos y otra informada por el retorno de mercado y el VIX, estableciendo las condiciones de equivalencia con el modelo nominal cuando el radio de robustez es cero.

\item \textbf{Evaluar} el desempeño fuera de muestra de las alternativas consideradas durante el periodo de prueba 2019–2025, manteniendo controlados los parámetros comunes para aislar el efecto de la geometría de transporte y del radio de robustez.

\item \textbf{Comparar} las alternativas consideradas en términos de crecimiento, pérdidas, concentración y costos de transacción, cuantificando además la incertidumbre asociada a las diferencias de crecimiento logarítmico.


%%% [C19b] Dir. p.8 (mapas de factibilidad, dffonsecav): con estas restricciones y δ<=0.90 el problema siempre es factible (ver Sec. de validación), por lo que el objetivo 5 ya no habla de «factibilidad» sino de estabilidad numérica.
%%% Texto anterior: \item \textbf{Analizar} la sensibilidad de la factibilidad y del desempeño del modelo Wasserstein–Kelly frente a cambios en el radio de robustez $\eta$ y en el parámetro de supervivencia $\delta$.
\item \textbf{Analizar} la sensibilidad del desempeño y de la estabilidad numérica del modelo Wasserstein–Kelly frente a cambios en el radio de robustez $\eta$ y en el parámetro de supervivencia $\delta$.
%%% [fin C19b]

\end{enumerate}

\subsection{Delimitación de la contribución}

Este trabajo no reclama ser el primero en proponer un portafolio de Kelly robusto
distribucionalmente, una bola de Wasserstein para retornos financieros, costos de
transacción en un portafolio logarítmico, o una reformulación convexa de un modelo WDRO
--todos estos elementos ya existen en la literatura citada arriba. La contribución es
específicamente \textbf{empírica y metodológica}: una comparación controlada, bajo un
protocolo de backtest idéntico, de si una geometría de transporte informada por el estado
del mercado aporta valor frente a una geometría más simple, dentro del contexto particular
de un portafolio Wasserstein--Kelly con costos de transacción.


% =====================================================================
% metodologia.tex -- BOCETO de la seccion de Metodologia
% Para usar: reemplaza el \section{Metodologia} vacio de main.tex por el
% contenido de este archivo (o usa \input{metodologia}).
%
% Convenciones:
%  - Lo marcado como \textbf{[PENDIENTE: ...]} NO esta hecho o decidido
%    todavia; buscalo con Ctrl+F "PENDIENTE" antes de entregar.
%  - Redactado en pasado (informe final). Las partes pendientes estan
%    escritas como se leeran cuando se completen.
%  - Solo usa los paquetes ya cargados en main.tex (amsmath, enumitem,
%    natbib, hyperref); no requiere booktabs.
%  - Al final del archivo hay entradas nuevas para thebibliography.
% =====================================================================

\section{Metodología}\label{sec:metodologia}

\subsection{Diseño general del estudio}\label{sec:diseno}

El estudio es una comparación computacional controlada de tres estrategias de asignación sobre el mismo universo de activos, el mismo periodo, las mismas restricciones y el mismo protocolo de backtest:
\begin{enumerate}[label=(\roman*)]
    \item el criterio de Kelly nominal \citep{kelly1956} con costos de transacción, que sirve de referencia;
    \item el modelo Wasserstein--Kelly con una geometría de transporte definida únicamente sobre los retornos de las acciones (geometría~1);
    \item el modelo Wasserstein--Kelly con una geometría informada además por el retorno agregado del mercado y el VIX (geometría~2).
\end{enumerate}

%%% [C5a] Dir. p.3 (identificación del efecto de la geometría): se elimina «difieren únicamente» y se separan el contraste emparejado y el de estrategias seleccionadas; Liu et al. solo justifica estudiar la geometría.
%%% Texto anterior: La comparación entre (ii) y (iii), que difieren únicamente en la geometría de transporte \citep{liu2022}, responde directamente a la pregunta de investigación; (i) permite medir qué aporta la robustez distribucional en sí misma.
La comparación entre (ii) y (iii), que se diferencian en la geometría de transporte, responde a la pregunta de investigación. Solo el contraste con parámetros $(\delta,\eta)$ emparejados aísla el efecto de la geometría; el contraste entre las configuraciones seleccionadas por validación para cada geometría compara, en cambio, estrategias calibradas, en las que cambian a la vez la geometría y los parámetros (Sección~\ref{sec:contrastes}). La cita de \citet{liu2022} respalda estudiar la geometría en DRO de manera general, pero la identificación experimental de su efecto proviene del contraste emparejado. La estrategia (i) permite medir qué aporta la robustez distribucional en sí misma.
%%% [fin C5a]
 En cada fecha de rebalanceo se resuelve un problema de optimización convexa con información disponible hasta ese momento y la decisión se evalúa con el retorno del mes siguiente (backtest de origen móvil). La Tabla~\ref{tab:objetivos} relaciona cada objetivo específico con la subsección que lo desarrolla.

\begin{table}[h]
\centering
\caption{Correspondencia entre objetivos específicos y metodología.}\label{tab:objetivos}
\begin{tabular}{|p{6.3cm}|p{8.7cm}|}
\hline
\textbf{Objetivo específico} & \textbf{Subsecciones} \\ \hline
1. Desempeño de referencia del Kelly nominal & \ref{sec:decision} (modelo de decisión y Kelly nominal) \\ \hline
2. Formular y validar el modelo Wasserstein--Kelly con dos geometrías & \ref{sec:wk} (modelo), \ref{sec:geom} (geometrías) y \ref{sec:impl} (verificación) \\ \hline
3. Evaluación fuera de muestra con parámetros comunes controlados & \ref{sec:protocolo} (partición temporal y protocolo) \\ \hline
4. Comparación de alternativas e incertidumbre de las diferencias & \ref{sec:metricas} (métricas e inferencia) \\ \hline
5. Sensibilidad a $\eta$ y a $\delta$ & \ref{sec:protocolo} (mallas de validación) \\ \hline
\end{tabular}
\end{table}

\subsection{Datos}\label{sec:datos}

\subsubsection{Universo y fuentes}

El universo invertible estuvo formado por $d=10$ acciones estadounidenses líquidas, 
%%% [C6] Dir. p.3 («de sectores distintos» no es correcto: AAPL/MSFT y PG/KO comparten sector). REVISAR: la clasificación GICS la propuse yo; verifícala con tu fuente.
%%% Texto anterior: de sectores distintos y con historia completa desde diciembre de 2009: AAPL, MSFT, AMZN, JPM, JNJ, XOM, PG, CAT, KO y WMT.
pertenecientes a varios sectores (tecnología de la información: AAPL y MSFT; consumo discrecional: AMZN; finanzas: JPM; salud: JNJ; energía: XOM; consumo básico: PG, KO y WMT; industriales: CAT), de modo que algunos sectores están representados por más de una acción, y con historia completa desde diciembre de 2009: AAPL, MSFT, AMZN, JPM, JNJ, XOM, PG, CAT, KO y WMT. Los criterios de selección fueron la liquidez y la disponibilidad de historia completa; no se buscó un equilibrio sectorial. 
%%% [fin C6]
 Se utilizaron precios diarios ajustados de Yahoo Finance\footnote{\url{https://finance.yahoo.com}}, el ETF SPY como aproximación al mercado accionario estadounidense, el índice VIX de volatilidad implícita \citep{whaley2000}\footnote{\url{https://www.cboe.com/tradable_products/vix/}}, y la tasa del Tesoro a tres meses (serie DGS3MO de FRED\footnote{\url{https://fred.stlouisfed.org/series/DGS3MO}}) como aproximación del activo de bajo riesgo. 
%%% [C9] Dir. p.4 (SPY como «benchmark externo»): se elimina esa promesa. Alternativa: incluir explícitamente una estrategia SPY comprar-y-mantener con su supuesto de costos, si se decide hacer esa comparación.
%%% Texto anterior: SPY no es una variable de decisión: se usó solo como benchmark externo y para calcular el retorno agregado de mercado.
SPY no es una variable de decisión: se usó únicamente como serie para calcular el retorno agregado de mercado, y no se incluyó como estrategia de comparación.
%%% [fin C9]
 Los archivos originales se conservaron sin modificación y se auditaron fechas duplicadas o faltantes, precios no positivos, retornos extremos y 
%%% [C7] Dir. p.4 (reproducibilidad de la auditoría): se agrega el tipo de cierre y los datos de descarga. Lo que depende de tus archivos de auditoría queda como PENDIENTE.
%%% Texto anterior: divisiones de acciones contra los datos históricos extraídos de Stooq\footnote{\url{https://stooq.com}}.

%%% [C7b] Auditoría (audit_report.md, sec. 1, 6 y 8): se reemplazan los PENDIENTE por los datos reales. Stooq bloqueado -> contraste con Tiingo/CBOE. REVISAR: «los 10 casos son de dic-2025» se deduce de mes_err_max=2025-12-31 en cada serie; confírmalo en audit_discrepancies.csv. NOTA: dic-2025 usa el cierre del 30-dic porque yf.download trata end como exclusivo; decide si regeneras (end=2026-01-01), lo que obliga a repetir los backtests.
%%% Texto anterior: divisiones de acciones contra los datos históricos extraídos de Stooq\footnote{\url{https://stooq.com}}. Se usaron cierres ajustados por dividendos y desdoblamientos (\emph{Adj Close} de Yahoo Finance), y la tasa DGS3MO se completó con su última observación disponible cuando faltaba un dato. Los datos se descargaron el \textbf{[PENDIENTE: fecha de descarga, registrada en \texttt{download\_log.csv}]}; el periodo efectivo por ticker fue \textbf{[PENDIENTE: confirmar con \texttt{audit\_report.csv}]}. Ante datos faltantes se \textbf{[PENDIENTE: criterio usado]}, y cuando Yahoo Finance y Stooq discrepaban \textbf{[PENDIENTE: regla de resolución; p.~ej., prevaleció Yahoo Finance y se registró la discrepancia]}.
divisiones de acciones, contrastando los precios con una fuente externa. Las fuentes usadas fueron Tiingo\footnote{\url{https://www.tiingo.com}} para las acciones y SPY, y contra los datos oficiales de CBOE\footnote{\url{https://www.cboe.com/tradable_products/vix/}} para el VIX (ambos descargados el 19 de octubre de 2026). Se usaron cierres ajustados por dividendos y desdoblamientos (\emph{Adj Close} de Yahoo Finance), y se verificó que el ajuste incorpora los dividendos y los cinco desdoblamientos del periodo (AAPL en 2014 y 2020, AMZN en 2022, KO en 2012 y WMT en 2024). Como Yahoo Finance reescala retroactivamente el \emph{Adj Close} con cada nuevo dividendo, los niveles dependen de la fecha de descarga, pero los retornos no.

Los datos de Yahoo Finance se descargaron el 12 de agosto de 2026; la serie DGS3MO se guardó junto con ellos, pero la fecha exacta de su descarga no quedó registrada. Las 12 series diarias de Yahoo cubren de diciembre de 2009 al 30 de diciembre de 2025 (4045 días hábiles): se solicitó hasta el 31 de diciembre, pero el parámetro de fin de la descarga es exclusivo, de modo que el cierre de diciembre de 2025 de las acciones, SPY y el VIX corresponde al 30 de diciembre. El panel mensual abarca de enero de 2010 a diciembre de 2025 (192 meses) sin valores faltantes. No se imputó ningún dato: el valor de cada mes es el último dato disponible del mes, nunca uno posterior; el único día hábil ausente (31 de diciembre de 2025) no se rellenó, y los NaN de DGS3MO, que corresponden a feriados de FRED, se omitieron tomando el último rendimiento válido del mes (ningún mes requirió completarse con datos de un mes anterior). No hubo fechas duplicadas ni precios no positivos, y el único retorno mensual con $|r|>30\,\%$ (CAT, octubre de 2022, $+32.8\,\%$) corresponde a un movimiento real posterior a sus resultados del tercer trimestre, no a un error de ajuste.

En el contraste, una discrepancia se definió como un error relativo de precio a fin de mes superior a $0.5\,\%$ o una diferencia de retorno total mensual superior a $0.25$ puntos porcentuales. Sobre las 12 series (193 meses cada una), 14 casos activo-mes superaron la tolerancia: 10 se explican porque el último día con dato del mes difiere entre fuentes (todos asociados al cierre de diciembre de 2025) y 4 quedaron sin explicar (AAPL en septiembre y octubre de 2015, y el VIX en agosto y septiembre de 2014); la mayor diferencia de retorno total, $3.8$ puntos porcentuales, ocurrió en el VIX. Se mantuvo Yahoo Finance como fuente oficial: las discrepancias se registraron y se clasificaron por causa probable, sin modificar ningún dato. 
%%% [fin C7b]

%%% [fin C7]


\subsubsection{Retornos mensuales y activo de bajo riesgo}

Con el último cierre ajustado de cada mes, $P_{j,m}$, el retorno simple mensual de la acción $j$ fue
\begin{equation}\label{eq:ret}
r_{j,m}=\frac{P_{j,m}}{P_{j,m-1}}-1,
\end{equation}
y el retorno mensual de mercado $m_m$ se calculó igual con SPY. Para cada mes se tomó el último valor diario del VIX, $V_m$. Si $y_k$ es el rendimiento anual porcentual de la serie DGS3MO al cierre del mes $k$, el retorno de bajo riesgo aplicable al periodo $k\to k+1$ fue
\begin{equation}\label{eq:rf}
r^f_k=\Bigl(1+\frac{y_k}{100}\Bigr)^{1/12}-1,
\end{equation}

%%% [C8] Dir. p.4 (disponibilidad temporal y alineación de fechas). REVISAR: describe lo que hace 01_import_data.py (resample('ME').last() + ffill); confirma el supuesto de observabilidad.
%%% Texto anterior: conocido en la fecha de decisión $k$. Esta transformación es un proxy
conocido en la fecha de decisión $k$. En cada mes se usaron el último cierre diario de SPY y del VIX y la última observación de DGS3MO publicada hasta el cierre del mes calendario; si la fecha de FRED no coincidía con el último día bursátil, 
%%% [C8b] Auditoría sec. 4, 5 y 8 (C8): se confirma la ausencia de look-ahead y se corrige la frase del ffill (el código no completó ningún mes con datos de uno anterior).
%%% Texto anterior: se tomó el último dato disponible hasta esa fecha (completando con el valor anterior cuando faltaba). Se supone que estas series son observables al cierre del mes $k$, instante en que se rebalancea. \textbf{[PENDIENTE: confirmar que ningún dato posterior al último día bursátil del mes interviene en la decisión.]}
se tomó el último dato disponible hasta esa fecha; los NaN de DGS3MO, que corresponden a feriados, se omitieron, y ningún mes requirió completarse con datos de un mes anterior. Se supone que estas series son observables al cierre del mes $k$, instante en que se rebalancea; en diciembre de 2025, SPY y el VIX usan el cierre del 30 de diciembre (el último dato descargado). Se verificó que ningún dato posterior al cierre de cada mes interviene en la decisión: al recomputar el panel desde los archivos originales no hubo celdas distintas (0 de 2496), y un contraste con retornos desplazados solo coincide con desfase cero.
%%% [fin C8b]
 Esta transformación es un proxy
%%% [fin C8]
 operativo del retorno de una posición de corto plazo, no el retorno realizado exacto de una cuenta monetaria.

\subsubsection{Partición temporal}

Los retornos mensuales cubren enero de 2010 en adelante. Los primeros 60 retornos (enero de 2010 a diciembre de 2014) formaron la primera ventana de estimación. El periodo de \emph{validación} (enero de 2015 a diciembre de 2018) se usó para depurar la implementación, explorar mallas de parámetros y elegir configuraciones; el periodo de \emph{prueba} (enero de 2019 a diciembre de 2025) se reservó para la evaluación final y no se usó para escoger ningún parámetro.

\subsection{Modelo de decisión en cada fecha de rebalanceo}\label{sec:decision}

En la fecha mensual $k$ el inversionista posee una riqueza $W_k>0$ distribuida según el portafolio heredado $h_k=(h_{0k},h_{1k},\dots,h_{dk})^\top$, con $h_{jk}\ge 0$ y $\sum_{j=0}^d h_{jk}=1$, donde el índice $0$ denota el activo de bajo riesgo. Sean $b_{jk}\ge0$ y $s_{jk}\ge0$ las compras y ventas de la acción $j$, expresadas como fracción de $W_k$, y $\tau\ge0$ el costo proporcional por unidad monetaria negociada \citep{hsiehyu2024}. Las posiciones después de operar y de pagar los costos son
\begin{align}
x_{jk}&=h_{jk}+b_{jk}-s_{jk}, \qquad j=1,\dots,d, \label{eq:bal_j}\\
x_{0k}&=h_{0k}-\sum_{j=1}^{d}b_{jk}+\sum_{j=1}^{d}s_{jk}-\tau\sum_{j=1}^{d}(b_{jk}+s_{jk}). \label{eq:bal_0}
\end{align}
Con $\mathrm{TO}_k=\sum_{j=1}^d(b_{jk}+s_{jk})$, esto implica $x_{0k}+\sum_{j=1}^d x_{jk}=1-\tau\,\mathrm{TO}_k$: las posiciones se expresan respecto a $W_k$ y la diferencia con uno es exactamente la fracción de riqueza consumida por los costos. Las compras y ventas simultáneas del mismo activo no son óptimas cuando $\tau>0$, por lo que no se requieren variables binarias.

En la fecha $k$ se dispone de la ventana de los $N$ retornos mensuales más recientes, $\mathcal D_k=\{\hat r_{k1},\dots,\hat r_{kN}\}$ con $\hat r_{ks}=r_{k-N+1+(s-1)}$, 
%%% [C10] Dir. p.5 (equiprobabilidad): se aclara que es la distribución empírica de referencia, no un supuesto i.i.d./estacionario.
%%% Texto anterior: que se interpretan como $N$ escenarios equiprobables para evaluar un portafolio candidato.
que se interpretan como $N$ escenarios equiprobables para evaluar un portafolio candidato. La equiprobabilidad es la distribución empírica de referencia de la ventana; no supone que los retornos sean independientes, estacionarios ni idénticamente distribuidos, y esta limitación se retoma en la Sección~\ref{sec:limit}.
%%% [fin C10]
 El factor de crecimiento del portafolio $x_k$ bajo el escenario $s$ es
\begin{equation}\label{eq:G}
G_{ks}(x_k)=x_{0k}(1+r^f_k)+\sum_{j=1}^{d}x_{jk}\,(1+\hat r_{jks}),
\end{equation}
una función afín de $x_k$ que, al incorporar $x_{0k}$, ya descuenta los costos de transacción.

\subsubsection{Kelly nominal con costos de transacción}

El modelo nominal maximiza el crecimiento logarítmico promedio sobre los escenarios de la ventana:
\begin{equation}\label{eq:nominal}
\begin{aligned}
\max_{x_k,b_k,s_k}\quad & \frac1N\sum_{s=1}^N\log G_{ks}(x_k)\\
\text{s.a.}\quad & \eqref{eq:bal_j},\ \eqref{eq:bal_0},\\
& 0\le x_{jk}\le \bar x,\quad \sum_{j=1}^d x_{jk}\le\gamma,\\
& G_{ks}(x_k)\ge\delta,\quad s=1,\dots,N,\\
& x_{0k}\ge0,\ b_{jk}\ge0,\ s_{jk}\ge0,
\end{aligned}
\end{equation}
donde $\bar x$ limita la posición por acción, $\gamma$ limita la exposición riesgosa total y $\delta>0$ es una cota de supervivencia \citep{hsiehyu2024} que exige que, bajo cada escenario de la ventana, la riqueza no caiga por debajo de la fracción $\delta$ de la riqueza previa a operar (además de garantizar que el logaritmo esté bien definido). El problema es convexo: maximiza una función cóncava sobre un conjunto convexo \citep{kelly1956,breiman1961}. La decisión $x_k$ se mantiene únicamente durante el periodo $k\to k+1$ y el problema se resuelve de nuevo en la fecha siguiente (horizonte rodante).

\subsection{Modelo Wasserstein--Kelly}\label{sec:wk}

El modelo nominal trata la distribución empírica de la ventana como si fuera la verdadera. Para protegerse contra errores en esa estimación \citep{li2023,sunboyd2018,hsiehyu2024}, se sustituye el promedio empírico por el peor caso dentro de una bola de Wasserstein de radio $\rho_k$ centrada en la distribución empírica, 
%%% [C11] Dir. p.5 (soporte restringido): se explica en una oración que es una bola discreta sobre el soporte empírico y se justifica por tratabilidad y alcance.
%%% Texto anterior: con soporte restringido a los $N$ escenarios históricos.
con soporte restringido a los $N$ escenarios históricos. Se trata, por tanto, de una bola de Wasserstein discreta restringida al soporte empírico: el adversario solo redistribuye probabilidad entre meses ya observados y no genera pérdidas fuera de ese soporte, de modo que las garantías de trabajos con soporte no restringido no se trasladan automáticamente a este caso. La restricción se adopta por tratabilidad (el dual \eqref{eq:dual} es finito y tiene $N^2$ restricciones) y por el alcance de la comparación, que contrasta geometrías sobre los mismos escenarios.
%%% [fin C11]
 Una distribución alternativa se describe por probabilidades $p=(p_1,\dots,p_N)$ en el símplex $\Delta_N$, y su distancia a la empírica $\hat p=(1/N)\mathbf 1$ es el costo mínimo de un plan de transporte \citep{mohajerin2018,kuhn2019}:
\begin{equation}\label{eq:W}
W_{C_k}(\hat p,p)=\min_{\Gamma\ge0}\Bigl\{\sum_{r,s}c_{k,rs}\Gamma_{rs}\;:\;\sum_{s}\Gamma_{rs}=\tfrac1N\ \forall r,\ \ \sum_{r}\Gamma_{rs}=p_s\ \forall s\Bigr\},
\end{equation}
donde $c_{k,rs}\ge0$ es el costo de trasladar probabilidad del escenario $r$ al escenario $s$ (definido en la Sección~\ref{sec:geom}). El conjunto de ambigüedad es $\mathcal P_{\rho_k}=\{p\in\Delta_N: W_{C_k}(\hat p,p)\le\rho_k\}$ y el problema de decisión es
\begin{equation}\label{eq:minimax}
\max_{x_k,b_k,s_k}\ \min_{p\in\mathcal P_{\rho_k}}\ \sum_{s=1}^N p_s\log G_{ks}(x_k)\qquad\text{sujeto a las mismas restricciones de \eqref{eq:nominal}}.
\end{equation}
El radio se normalizó con la escala de distancias de cada ventana,
\begin{equation}\label{eq:rho}
\rho_k=\eta\cdot\operatorname{mediana}\{c_{k,rs}: r<s,\ c_{k,rs}>0\},
\end{equation}

%%% [C12/C13a] Dir. p.6 (comparabilidad de η y dualización). REVISAR: la regla «sin distancias positivas => ρ_k=0» debe coincidir con lo que hace compute_rho.m.
%%% Texto anterior: de modo que $\eta\ge0$ es una intensidad \emph{relativa} de robustez comparable entre ventanas y entre geometrías. Por dualidad lineal del problema de transporte \citep{mohajerin2018,kuhn2019}, el problema \eqref{eq:minimax} es equivalente a la reformulación convexa exacta
de modo que $\eta\ge0$ es una intensidad \emph{relativa} que normaliza la escala de los costos en cada ventana. Esta normalización no implica que un mismo $\eta$ represente exactamente la misma severidad de ambigüedad bajo geometrías distintas, pues sus distribuciones de distancias y los conjuntos de probabilidades alcanzables pueden diferir; por eso la comparación entre geometrías se hace sobre una malla de $\eta$ y no se interpreta como equivalencia estricta de radios. Si en una ventana no existe ningún costo positivo fuera de la diagonal, se toma $\rho_k=0$ y el modelo se reduce al nominal. 
%%% [R12] Reporte T2-T5 (resumen_distancias_c.csv; compute_rho.m): se reemplaza el PENDIENTE opcional por el resumen de distancias. La comparación del objetivo en el peor caso a un eta común NO se hizo y se declara así.
%%% Texto anterior: [PENDIENTE (opcional): resumir (mediana y cuartiles) las distancias c_{k,rs} de cada geometría y, si se dispone, comparar el valor del objetivo en el peor caso a un eta común.]
Se resumieron (mediana y cuartiles) los costos $c_{k,rs}$ de cada ventana para cada geometría; no se comparó el valor del objetivo en el peor caso a un $\eta$ común.
%%% [fin R12]
Con soporte finito, el problema interno de \eqref{eq:minimax} es un programa lineal en el plan de transporte, y su dualidad lineal \citep[como en][]{mohajerin2018,kuhn2019} conduce a la reformulación convexa exacta
%%% [fin C12/C13a]

\begin{equation}\label{eq:dual}
\begin{aligned}
\max_{x_k,b_k,s_k,\alpha,\lambda}\quad & \frac1N\sum_{r=1}^N\alpha_r-\rho_k\lambda\\
\text{s.a.}\quad & \alpha_r\le\log G_{ks}(x_k)+\lambda\,c_{k,rs},\quad r,s=1,\dots,N,\\
& \lambda\ge0,\ \alpha\in\mathbb R^N,\ \text{y las restricciones de \eqref{eq:nominal}},
\end{aligned}
\end{equation}

%%% [C13b] Dir. p.6 (justificación de la dualidad): se agrega una derivación breve, el papel de λ y α y la hipótesis de positividad dada por δ.
%%% Texto anterior: que es representable con el cono exponencial y tiene $N^2$ restricciones logarítmicas. Con $\rho_k=0$ y costos fuera de la diagonal estrictamente positivos se recupera exactamente el modelo nominal \eqref{eq:nominal}.
que es representable con el cono exponencial y tiene $N^2$ restricciones logarítmicas. Con $\rho_k=0$ y costos fuera de la diagonal estrictamente positivos se recupera exactamente el modelo nominal \eqref{eq:nominal}.

\emph{Justificación de la dualización.} Fijado $x_k$, sea $\ell_s=\log G_{ks}(x_k)$. Dualizando con $\lambda\ge0$ la restricción $\sum_{r,s}c_{k,rs}\Gamma_{rs}\le\rho_k$, el problema interno vale $\max_{\lambda\ge0}\bigl\{-\rho_k\lambda+\tfrac1N\sum_{r}\min_{s}(\ell_s+\lambda c_{k,rs})\bigr\}$; la variable $\alpha_r$ representa $\min_s(\ell_s+\lambda c_{k,rs})$ y se impone mediante las $N^2$ desigualdades de \eqref{eq:dual}. Así, $\lambda$ es el precio de la restricción de presupuesto de transporte y $\alpha_r$ el valor del peor destino disponible para la masa del escenario $r$. La ausencia de brecha de dualidad proviene de la dualidad lineal de este problema discreto, cuyo conjunto factible es no vacío porque contiene a $\hat p$; la positividad de $G_{ks}$, garantizada por $\delta>0$, asegura que $\ell_s$ esté bien definido, y la concavidad de $\log G_{ks}$ en $x_k$ hace convexas las restricciones. La exactitud de \eqref{eq:dual} se deriva, por tanto, de esta estructura discreta; \citet{mohajerin2018} y \citet{kuhn2019} se citan como antecedentes de la dualización de Wasserstein, no como fuente de esta fórmula para este soporte y esta función logarítmica.
%%% [fin C13b]


\subsection{Geometrías de transporte}\label{sec:geom}

Los costos $c_{k,rs}$ se construyen en cada fecha usando únicamente la ventana $\mathcal D_k$. Cada variable se estandariza con la mediana y el rango intercuartílico de la ventana,
\begin{equation}\label{eq:std}
\tilde z_{ks}=\frac{z_{ks}-\operatorname{mediana}\{z_{k1},\dots,z_{kN}\}}{\max\{\operatorname{IQR}(z_{k1},\dots,z_{kN}),\,\varepsilon_{\mathrm{esc}}\}},
\end{equation}

%%% [C14] Dir. p.6 (piso ε_esc). REVISAR: la regla «excluir la variable constante» es la sugerida por el director y no está implementada; ajusta si el conteo es cero.
%%% Texto anterior: para que ninguna variable domine la distancia por su escala numérica. Los valores originales de los retornos siguen apareciendo en \eqref{eq:G}; la estandarización solo define la noción de cercanía entre escenarios.
para que ninguna variable domine la distancia por su escala numérica. Los valores originales de los retornos siguen apareciendo en \eqref{eq:G}; la estandarización solo define la noción de cercanía entre escenarios. El piso $\varepsilon_{\mathrm{esc}}$ evita divisiones por cero, pero puede amplificar variaciones mínimas si una serie tiene un IQR casi nulo en la ventana (en particular una covariable). 
%%% [C14b] Auditoría sec. 7 y 8 (C14): conteo = 0 en 132 ventanas x 12 series.
%%% Texto anterior: Por ello se contó el número de ventanas con $\mathrm{IQR}<\varepsilon_{\mathrm{esc}}$ en cada serie: \textbf{[PENDIENTE: reportar el conteo; si es cero, indicarlo]}. Si llegara a ocurrir, la regla prevista es excluir de $c_{k,rs}$ la contribución de la variable constante en esa ventana y renormalizar los pesos restantes.
Por ello se contó el número de ventanas con $\mathrm{IQR}<\varepsilon_{\mathrm{esc}}$: en las 132 ventanas de 60 meses del backtest y las 12 series de la geometría (10 retornos, SPY y VIX) el conteo fue cero; el IQR más pequeño ($0.026$) correspondió al retorno de SPY, muy por encima del piso. La regla de exclusión prevista (omitir de $c_{k,rs}$ la contribución de una variable constante en la ventana) no llegó a activarse.
%%% [fin C14b]

%%% [fin C14]


\paragraph{Geometría 1 (solo retornos).} El costo es la diferencia absoluta promedio por activo entre los retornos estandarizados:
\begin{equation}\label{eq:g1}
c^{R}_{k,rs}=\frac1d\,\bigl\|\tilde r_{kr}-\tilde r_{ks}\bigr\|_1 .
\end{equation}

\paragraph{Geometría 2 (retornos, mercado y VIX).} A cada escenario histórico se le asocia el retorno de mercado $m_{ks}$ y el nivel del VIX $v_{ks}$ del mes correspondiente, y el costo combina los tres bloques:
\begin{equation}\label{eq:g2}
c^{M}_{k,rs}=\omega_R\,c^{R}_{k,rs}+\omega_m\,|\tilde m_{kr}-\tilde m_{ks}|+\omega_v\,|\tilde v_{kr}-\tilde v_{ks}|,\qquad \omega_R+\omega_m+\omega_v=1,
\end{equation}

%%% [C15] Dir. p.7 (justificación de los pesos). REVISAR: confirma que esta es tu razón para 70/15/15; si es otra, reemplázala.
%%% Texto anterior: con $(\omega_R,\omega_m,\omega_v)=(0.70,\,0.15,\,0.15)$. La tasa de corto plazo
con $(\omega_R,\omega_m,\omega_v)=(0.70,\,0.15,\,0.15)$. Los pesos son una elección de diseño a priori, no calibrada con datos: se asignó $70\,\%$ a los retornos de los activos, que son la información directamente relevante para el portafolio, y $15\,\%$ a cada variable de estado, de modo que las covariables modifiquen la noción de cercanía sin dominarla. Como al cambiar los pesos varía también el aporte relativo de los retornos, las diferencias observadas entre geometrías no se atribuyen solo a «incorporar información», y las conclusiones están condicionadas a esta elección.
%\textbf{[PENDIENTE (opcional): comprobación con una distribución alternativa de pesos, sin optimizarla sobre el periodo de prueba.]} La tasa de corto plazo
%%% [fin C15]
 se usó únicamente como $r^f_k$ en \eqref{eq:G}, no como componente de la distancia. Dos meses con retornos accionarios parecidos pero con niveles de VIX muy distintos quedan así más lejos que bajo la geometría~1. 
%%% [C16] Dir. p.7 (la geometría no condiciona la distribución al estado actual) y p.2 (motivación, no evidencia de mejora).
%%% Texto anterior: Las variables de estado no aparecen en el factor de crecimiento ni se usan para predecir el régimen siguiente: solo impactan la ponderación a través de la distancia. La elección del VIX y del retorno de mercado como variables de estado se justificó en la Sección~\ref{sec:vacio} \citep{whaley2000,french1987}.
Las variables de estado no aparecen en el factor de crecimiento ni se usan para predecir el régimen siguiente: solo intervienen en la distancia entre escenarios. El estado actual del mercado no condiciona explícitamente la distribución predictiva; solo la ventana histórica y su geometría determinan el conjunto de ambigüedad. La geometría usa los estados históricos para definir los costos \emph{entre} escenarios y no aproxima una distribución condicional al VIX o al retorno de mercado observados hoy, por lo que el modelo no reconoce ni anticipa regímenes. La elección del VIX y del retorno de mercado se motivó en la Sección~\ref{sec:vacio} \citep{whaley2000,french1987} como una hipótesis económicamente razonable que este estudio pone a prueba, no como un resultado supuesto.
%%% [fin C16]


\subsection{Protocolo experimental}\label{sec:protocolo}

\subsubsection{Backtest de origen móvil}


%%% [C17] Dir. p.7 (condición inicial por segmento). REVISAR: se redactó según extract_period_results.m (recorte de la trayectoria completa); confirma antes de enviar.
%%% Texto anterior: Se partió con $W_0=1$ y todo el capital en el activo de bajo riesgo, $h_0=(1,0,\dots,0)^\top$, de modo que la primera asignación también pagó costos.
Se partió con $W_0=1$ y todo el capital en el activo de bajo riesgo, $h_0=(1,0,\dots,0)^\top$, de modo que la primera asignación también pagó costos. Esta condición inicial se aplicó una sola vez por trayectoria, en la primera fecha de decisión (diciembre de 2014), para cada estrategia y para cada combinación de la malla. El tramo de prueba no reinició posiciones: sus resultados provienen de la misma trayectoria, por lo que parte del portafolio heredado $h_k$ acumulado al cierre de 2018, y la riqueza se re-basó a uno al inicio del tramo. La misma regla se aplicó a todas las estrategias comparadas. No se corrió un backtest alternativo que empiece en efectivo en 2019.
%%% [R17b] Resumen T10 (T8): el portafolio heredado al 31-dic-2018 tenia h_0=0 y suma de h_j=1 en las 6 estrategias; no se corrio un backtest desde efectivo en 2019. Se agrega como precision opcional.
%%% [R17] Reporte T2-T5 (main_test_evaluation.m y main_test_evaluation_v2.m): confirmado que el tramo de prueba se obtiene recortando la trayectoria completa. Se reemplaza el PENDIENTE.
%%% Texto anterior: [PENDIENTE: confirmar que así se ejecutó la evaluación de prueba, es decir, recortando la trayectoria completa y no corriendo un backtest nuevo desde efectivo en 2019.]
La evaluación de prueba se ejecutó así: se recortó la trayectoria completa guardada, que comienza en diciembre de 2014 con todo el capital en efectivo, y no se corrió un backtest nuevo desde efectivo en 2019.
%%% [fin R17]
%%% [fin C17]
 En cada fecha mensual $k$ se siguió el mismo procedimiento para los tres modelos:
\begin{enumerate}[label=\arabic*.]
    \item se tomó la información disponible hasta el cierre del mes $k$ y se construyó la ventana $\mathcal D_k$ con $N=60$ meses;
    \item se estandarizó la ventana y se construyó $C_k$ (solo en los modelos Wasserstein--Kelly) y $\rho_k$ mediante \eqref{eq:rho};
    \item se resolvió el problema correspondiente --\eqref{eq:nominal} o \eqref{eq:dual}-- y se obtuvo $x_k$;
    \item se mantuvo $x_k$ durante el mes siguiente y se observó el retorno realizado $r_{k+1}$, con factor de crecimiento
    \begin{equation}\label{eq:Greal}
    G^{\mathrm{real}}_{k+1}=x_{0k}(1+r^f_k)+\sum_{j=1}^d x_{jk}(1+r_{j,k+1}),\qquad W_{k+1}=W_k\,G^{\mathrm{real}}_{k+1};
    \end{equation}
    \item se actualizó el portafolio heredado, $h_{0,k+1}=x_{0k}(1+r^f_k)/G^{\mathrm{real}}_{k+1}$ y $h_{j,k+1}=x_{jk}(1+r_{j,k+1})/G^{\mathrm{real}}_{k+1}$, se desplazó la ventana y se repitió.
\end{enumerate}
Ninguna observación posterior a $k$ interviene en $x_k$. Los costos no se descuentan de nuevo en \eqref{eq:Greal} porque ya están incorporados en $x_{0k}$. 
%%% [C18] Dir. p.7 (fallo del solucionador != infactibilidad). REVISAR: confirma que los estados exactos de MOSEK se guardan por mes.
%%% Texto anterior: Si en alguna fecha el solucionador no devolvía una solución óptima, el portafolio heredado se mantuvo sin rebalancear ese mes y el evento se registró como fecha no factible.
En cada fecha se registró el estado devuelto por el solucionador. Que MOSEK/CVX no declare una solución óptima no implica que el modelo matemático sea infactible: puede deberse a un límite de tiempo, a inestabilidad numérica o a un estado \emph{inaccurate}. Por ello se distinguió entre infactibilidad certificada y fallo numérico o no resuelto, y se aplicó la misma regla de contingencia a todas las estrategias: se mantuvo el portafolio heredado ese mes y el evento se contó como \emph{mes sin optimización}, no automáticamente como problema infactible, porque esa decisión alternativa altera retornos y costos.
%%% [R18] Reporte T2-T5, sec. 2: precisión de CVX/MOSEK, estados guardados y conteo de meses sin optimización. Se reemplaza el PENDIENTE.
%%% Texto anterior: [PENDIENTE: indicar la precisión del solucionador y confirmar que el estado exacto de cada mes quedó guardado.]
En cada mes se guardaron el estado devuelto por CVX, el estado nativo de MOSEK (\emph{Problem status} y \emph{Solution status}) y una clasificación en \{óptimo, \emph{inaccurate}, infactible certificado, fallo numérico\}; el estado nativo de MOSEK se guarda solo en las corridas hechas después de incorporar este registro, mientras que los resultados guardados antes conservan únicamente el estado de CVX, de donde se derivó su clase. La precisión fue la predeterminada de CVX 2.2, $\texttt{cvx\_precision}=[1.49\times10^{-8},\,1.49\times10^{-8},\,1.22\times10^{-4}]$.
%%% [fin R18]
%%% [fin C18]
 La cota $G_{ks}(x_k)\ge\delta$ se impone sobre los escenarios de la ventana, por lo que no garantiza que el factor realizado sea mayor que $\delta$; cada incumplimiento fuera de muestra se registró.

\subsubsection{Parámetros comunes y mallas}

Los parámetros que se mantuvieron fijos entre los tres modelos para aislar el efecto de la geometría y del radio se resumen en la Tabla~\ref{tab:param}. Solo $\delta$ y $\eta$ se variaron, en las mallas de validación.

\begin{table}[h]
\centering
\caption{Parámetros comunes y mallas de validación.}\label{tab:param}
\begin{tabular}{|p{6.0cm}|p{9.0cm}|}
\hline
\textbf{Parámetro} & \textbf{Valor} \\ \hline
Ventana de estimación $N$ & 60 meses \\ \hline
Costo proporcional $\tau$ & $0.001$ (10 puntos básicos por unidad negociada) \\ \hline
Límite por acción $\bar x$ / exposición riesgosa $\gamma$ & $0.25$ / $1$ \\ \hline
Piso de estandarización $\varepsilon_{\mathrm{esc}}$ & $10^{-6}$ \\ \hline
Pesos de la geometría 2 $(\omega_R,\omega_m,\omega_v)$ & $(0.70,\,0.15,\,0.15)$ \\ \hline
Nivel de confianza de VaR y CVaR & $0.95$ \\ \hline
Malla de $\delta$ (validación) & $\{0.50,\,0.60,\,0.70,\,0.75,\,0.80,\,0.90\}$ \\ \hline

%%% [C26] Dir. p.8 (Cuadro 2 se salía del margen): la malla de η se parte en dos líneas.
%%% Texto anterior: Malla de $\eta$ (validación) & $\{0,\,0.01,\,0.02,\,0.03,\,0.04,\,0.05,\,0.10,\,0.15,\,0.30,\,0.50,\,1.00\}$ \\ \hline
Malla de $\eta$ (validación) & $\{0,\,0.01,\,0.02,\,0.03,\,0.04,\,0.05,$\newline $0.10,\,0.15,\,0.30,\,0.50,\,1.00\}$ \\ \hline
%%% [fin C26]

\end{tabular}
\end{table}

\subsubsection{Validación y prueba}

\emph{Fase de validación.} Para cada geometría se corrió la malla completa $(\delta,\eta)$ de la Tabla~\ref{tab:param} sobre el periodo de validación, usando datos recortados al final de 2018, de modo que el código no tuvo acceso al periodo de prueba. 
%%% [C19] Dir. p.8 (mapas de factibilidad, dffonsecav). Argumento verificado: x_0=1-τΣh_j>=0.999 y G=x_0(1+r_f)>0.90. REVISAR: si en tus mallas aparecen fracciones factibles <1, son fallos numéricos y deben reportarse como tales.
%%% Texto anterior: Para cada combinación se registraron la fracción de fechas factibles y las métricas de la Sección~\ref{sec:metricas}; con ellas se construyeron mapas de factibilidad y de desempeño sobre la malla $(\delta,\eta)$, que sustentan el análisis de sensibilidad del objetivo~5.
Para cada combinación se registraron la fracción de meses con solución óptima y las métricas de la Sección~\ref{sec:metricas}; con ellas se construyeron mapas de desempeño, de actividad de la cota $\delta$ y de fallos numéricos sobre la malla $(\delta,\eta)$, que sustentan el análisis de sensibilidad del objetivo~5. Con las restricciones del modelo y $\delta\le0.90$, el problema es siempre factible: vender todas las posiciones riesgosas ($x_{jk}=0$) deja $x_{0k}=1-\tau\sum_{j}h_{jk}\ge1-\tau=0.999$ y $G_{ks}=x_{0k}(1+r^f_k)\ge0.999\,(1+r^f_k)>0.90$ para las tasas mensuales observadas ($r^f_k\ge0$), de modo que la cota de supervivencia se cumple en todos los escenarios y el conjunto factible es no vacío para toda combinación $(\delta,\eta)$ de la malla. Por esta razón no se buscaron regiones infactibles en sentido teórico: los mapas describen (a) las regiones donde la cota $\delta$ está activa, (b) el desempeño y el riesgo de pérdidas y (c) la frecuencia de meses sin optimización por fallos numéricos, que se reportan separados de la factibilidad teórica.
%%% [fin C19]
 El caso $\eta=0$ de cada fila de $\delta$ equivale al modelo nominal con ese mismo $\delta$ y se usó como línea base.
%%% [C20] Dir. p.8 (oración repetida): se elimina «De igual manera, se ejecutó el mismo protocolo ... geometría 2», que sugería una validación adicional.
%%% Texto anterior:  De igual manera, se ejecutó el mismo protocolo de validación para la geometría 2.
%%% [C20 | resuelto] la malla de validación de la geometría 2 ya estaba corrida (66 celdas, 48 meses, todas con solución óptima; reporte T2-T5, sec. 3), de modo que el párrafo anterior es cierto para ambas geometrías.
%%% [fin C20]


%%% [C21/C23] Dir. p.8 (regla de selección incompleta; criterio solo de riqueza). REVISAR: los pasos (2) y (3) —umbral φ_min y desempate— son propuestas mías; fíjalos antes de mirar la prueba.
%%% Texto anterior: \emph{Selección.} La configuración $(\delta,\eta)$ de cada geometría se eligió únicamente con los resultados de validación, aplicando una regla fijada de antemano. La ejecución de la malla El criterio a priori de elección de dicha configuración fue el cociente entre riqueza final e inicial. A partir de la elección de parámetros en ambas geometrías, y se comparan tanto los mejores modelos resultantes de cada geometría, como sus mismas versiones en ambas geometrías. 
%%% [R21] Reportes T9 (REPORTE_SELECCION.md): se fijan phi_min, TOL y el paso 1 (solo eta>0). Los valores elegidos y su sensibilidad son resultados y no se incluyen aqui. Texto anterior (comentado):
%%%   \emph{Selección.} La configuración $(\delta,\eta)$ de cada geometría se eligió únicamente con los resultados de validación, mediante una regla fijada antes de observar el periodo de prueba y compuesta de tres pasos: (1) la métrica primaria es el cociente entre la riqueza final y la inicial en validación, que se maximiza; (2) una configuración es elegible si su fracción de meses con solución óptima es al menos $\phi_{\min}$ 
%%%   %\textbf{[PENDIENTE: fijar $\phi_{\min}$ antes de usarlo]} 
%%%   y sus fallos se trataron con la misma regla de contingencia que las demás; y (3) los empates (diferencias relativas de riqueza inferiores a $10^{-6}$) se rompen eligiendo el menor $\eta$ (la configuración más cercana al modelo nominal) y, si persiste el empate, el mayor $\delta$. La riqueza es el criterio primario; crecimiento, pérdidas, concentración y costos son dimensiones secundarias que se reportan, pero no se usaron para elegir a posteriori, por lo que la configuración elegida no es la mejor en un sentido multicriterio. Con 48 meses de validación y 66 combinaciones $(\delta,\eta)$ por geometría existe riesgo de sobreajustar la validación, lo que se reconoce como limitación (Sección~\ref{sec:limit}).
\emph{Selección.} La configuración $(\delta,\eta)$ de cada geometría se eligió únicamente con los resultados de validación (48 meses realizados, con decisiones de diciembre de 2014 a noviembre de 2018), mediante una regla aplicada por separado a cada geometría y compuesta de cuatro pasos: (1) solo son elegibles las celdas con $\eta>0$, porque $\eta=0$ es el Kelly nominal, que se mantiene como referencia aparte y en el cual las dos geometrías coinciden; (2) una celda es elegible si su fracción de meses con solución óptima es al menos $\phi_{\min}=0.95$; (3) entre las elegibles se maximiza el cociente entre la riqueza final y la inicial en validación; y (4) las celdas cuya riqueza difiere de la máxima en menos de una tolerancia relativa $\mathrm{TOL}=5\times10^{-4}$ se consideran empatadas, y el empate se rompe eligiendo el menor $\eta$ y, si persiste, el mayor $\delta$. La tolerancia se fijó a partir del ruido numérico medido en validación: se registró en cada celda el número de meses con la cota de supervivencia activa (holgura $\le10^{-6}$); las celdas con un mismo $\eta$ en las que la cota nunca está activa son la misma estrategia en teoría, y la diferencia relativa máxima de riqueza final entre ellas se tomó como ruido numérico. $\mathrm{TOL}$ es ese máximo redondeado hacia arriba a la siguiente cifra de la forma $\{1,2,5\}\times10^{k}$. La exclusión de $\eta=0$ (paso 1) se decidió después de observar que, con la regla original, el nominal ocupaba los primeros puestos de la tabla de validación; por eso la regla no se fijó antes de ver esa tabla. La riqueza es el criterio primario; crecimiento, pérdidas, concentración y costos son dimensiones secundarias que se reportan, pero no se usaron para elegir a posteriori, por lo que la configuración elegida no es la mejor en un sentido multicriterio. Con 48 meses de validación y 60 celdas elegibles por geometría existe riesgo de sobreajustar la validación, lo que se reconoce como limitación (Sección~\ref{sec:limit}). Se denotan $(\delta_1^\star,\eta_1^\star)$ y $(\delta_2^\star,\eta_2^\star)$ las configuraciones elegidas para las geometrías~1 y 2.
%%% [fin R21]
%%% [fin C21/C23]

\emph{Fase de prueba.} 
%%% [C22/C24] Dir. p.8 (contraste emparejado vs. seleccionado) y p.8 (la equivalencia «recortar = correr por separado» no es correcta con costos y portafolio heredado): se reemplaza por la redacción sugerida.
%%% Texto anterior: Las configuraciones elegidas se evaluaron una sola vez sobre el periodo de prueba, con los parámetros congelados, junto con el modelo nominal y comparando cada alternativa contra las demás. Los resultados de cada tramo se obtienen recortando por fecha la misma trayectoria del backtest, y la riqueza se re-basa a uno al inicio del tramo; esto es equivalente a correr cada tramo por separado, porque cada decisión solo usa información pasada.
%%% [R22] Reporte de prueba (REPORTE_PRUEBA.md): se detalla qué se corrió. Se retira «una sola vez» porque la configuración de la geometría 1 ya se había evaluado en prueba antes de T9 (ver limitaciones). Texto anterior:
%%% Las configuraciones elegidas se evaluaron una sola vez sobre el periodo de prueba, con los parámetros congelados, junto con el modelo nominal. Los resultados por tramo se obtuvieron recortando la trayectoria ya simulada y reescalando su riqueza al inicio del tramo, sin reiniciar posiciones ni reoptimizar.
Las configuraciones elegidas se evaluaron sobre el periodo de prueba con los parámetros congelados, junto con el modelo nominal. Como la malla de la geometría~2 solo se había corrido sobre validación, esta geometría se corrió sobre los 132 meses (decisiones de diciembre de 2014 a noviembre de 2025) únicamente para las configuraciones seleccionadas, y se verificó que sus primeros 48 meses coinciden con la trayectoria de validación guardada. Se evaluaron el Kelly nominal con el $\delta$ de cada configuración y las geometrías~1 y 2 con cada una de las dos configuraciones seleccionadas; la geometría~1 y el nominal provienen de las corridas ya guardadas. Los resultados por tramo se obtuvieron recortando la trayectoria ya simulada y reescalando su riqueza al inicio del tramo, sin reiniciar posiciones ni reoptimizar.
%%% [fin R22]

\paragraph{Contrastes de la fase de prueba.}\label{sec:contrastes} Se hicieron dos contrastes separados. (a) \emph{Contraste controlado:} las geometrías 1 y 2 se evaluaron con parámetros $(\delta,\eta)$ idénticos --cada una de las dos configuraciones seleccionadas, aplicada a ambas geometrías--, de modo que solo cambia la geometría; este contraste aísla su efecto. (b) \emph{Contraste entre estrategias seleccionadas:} cada geometría con su propia configuración $(\delta_g^\star,\eta_g^\star)$; compara estrategias calibradas en las que cambian a la vez la geometría y los parámetros, por lo que no identifica el efecto de la geometría. Cada geometría se contrastó además con el Kelly nominal con el mismo $\delta$ de la configuración. En el periodo de prueba no se usó ningún resultado para volver a elegir parámetros.
%%% [fin C22/C24]
 No se barrió de nuevo ninguna malla sobre el periodo de prueba.

\subsection{Métricas y evaluación estadística}\label{sec:metricas}

\subsubsection{Métricas de desempeño}

Para cada estrategia y cada tramo se calcularon, a partir de los factores de crecimiento realizados:
\begin{itemize}
    \item \textbf{Crecimiento:} riqueza final, crecimiento logarítmico medio y tasa de crecimiento anual compuesta,
    \begin{equation}\label{eq:metricas}
    g=\frac1T\sum_{t=1}^{T}\log\frac{W_t}{W_{t-1}},\qquad \mathrm{CAGR}=\Bigl(\frac{W_T}{W_0}\Bigr)^{12/T}-1 .
    \end{equation}
    \item \textbf{Pérdidas:} máximo drawdown, $\mathrm{MDD}=\max_t\bigl(1-W_t/\max_{u\le t}W_u\bigr)$, peor retorno mensual, varianza de los retornos mensuales, 
%%% [C25] Dir. p.8 (CVaR empírico). REVISAR: alinea compute_metrics.m con esta definición (hoy promedia las pérdidas >= VaR).
%%% Texto anterior: VaR y CVaR empíricos al $95\,\%$ de la pérdida mensual (el CVaR como promedio de las pérdidas que alcanzan o superan el cuantil \citep{rockafellaruryasev2000}).
VaR y CVaR empíricos al $95\,\%$ de la pérdida mensual $L_t=1-G^{\mathrm{real}}_t$ (el VaR es el cuantil de nivel $\alpha=0.95$ de $L_t$, tomado como el estadístico de orden $L_{(\lceil\alpha T\rceil)}$, que es el $\zeta$ óptimo de la formulación siguiente; el CVaR se calcula con la formulación empírica de \citet{rockafellaruryasev2000}, $\mathrm{CVaR}_\alpha=\min_{\zeta}\bigl\{\zeta+\tfrac{1}{(1-\alpha)T}\sum_{t}\max(L_t-\zeta,0)\bigr\}$, que maneja empates y colas de tamaño no entero). Con $T=84$ meses en prueba, la cola al $95\,\%$ contiene muy pocas observaciones, por lo que VaR y CVaR se presentan como medidas descriptivas de muestra con alta incertidumbre, no como estimaciones precisas de riesgo extremo.
%%% [R25] Reporte T2-T5 (sec. 1.1): compute_metrics.m define VaR = L_(ceil(alpha T)) y CVaR de Rockafellar-Uryasev exacto; se explicita el VaR como estadistico de orden.
%%% [fin C25]

    \item \textbf{Riesgo ajustado:} 
%%% [C27] Dir. p.9 (fórmulas, tasa objetivo del Sortino y denominador nulo). REVISAR: ajusta a lo que realmente calcula compute_metrics.m.
%%% Texto anterior: razones de Sharpe \citep{sharpe1994}, Sortino \citep{sortinoprice1994} (con desviación a la baja respecto a cero) y Calmar \citep{young1991} ($\mathrm{CAGR}/\mathrm{MDD}$), anualizadas.
razones de Sharpe \citep{sharpe1994}, Sortino \citep{sortinoprice1994} y Calmar \citep{young1991}. Con $R_t=G^{\mathrm{real}}_t-1$ y $e_t=R_t-r^f_{t-1}$, el Sharpe se calculó con excesos mensuales frente a $r^f$, $\sqrt{12}\,\bar e/s_e$; el Sortino con tasa objetivo cero (no $r^f$), $\sqrt{12}\,\bar R/\mathrm{DD}_0$ con $\mathrm{DD}_0=\bigl(\tfrac1T\sum_t\min(R_t,0)^2\bigr)^{1/2}$; y el Calmar como $\mathrm{CAGR}/\mathrm{MDD}$. Si el denominador es nulo, la razón se reportó como no definida.
%%% [fin C27]

    \item \textbf{Costos:} 
%%% [C28] Dir. p.9 (qué es «costo acumulado»): se distinguen el costo fraccional y el monetario.
%%% Texto anterior: turnover acumulado $\sum_k\mathrm{TO}_k$, número de meses con transacción, y costo acumulado y por rebalanceo.
turnover acumulado $\sum_k\mathrm{TO}_k$, número de meses con transacción, costo fraccional acumulado $\sum_k\tau\,\mathrm{TO}_k$ (intensidad de negociación; suma de fracciones de riquezas $W_k$ distintas) y costo monetario acumulado en unidades de riqueza inicial, $\sum_k W_k\,\tau\,\mathrm{TO}_k/W_0$, además del costo por rebalanceo.
%%% [fin C28]

    \item \textbf{Concentración:} 
%%% [C29] Dir. p.9 (N_eff indefinido sin exposición riesgosa). REVISAR: confirma que plot/compute usan N_eff=0 en ese caso.
%%% Texto anterior: peso máximo y número efectivo de activos riesgosos, $N_{\mathrm{eff},k}=1/\sum_{j}\pi_{jk}^2$ con $\pi_{jk}=x_{jk}/\sum_{\ell}x_{\ell k}$.
peso máximo y número efectivo de activos riesgosos, $N_{\mathrm{eff},k}=1/\sum_{j}\pi_{jk}^2$ con $\pi_{jk}=x_{jk}/\sum_{\ell}x_{\ell k}$. Cuando $\sum_\ell x_{\ell k}=0$ (portafolio totalmente en bajo riesgo), $\pi_{jk}$ no está definido: se fijó $N_{\mathrm{eff},k}=0$ y la concentración riesgosa se reportó como no aplicable. Como $N_{\mathrm{eff}}$ mide la diversificación entre activos riesgosos dado que hay exposición, y no la concentración del patrimonio total, se reporta por separado la proporción invertida en riesgo, $\sum_j x_{jk}$. En los promedios de $N_{\mathrm{eff}}$ y del peso máximo solo se incluyen los meses con exposición riesgosa ($\sum_\ell x_{\ell k}>10^{-8}$), y se reporta el número de meses sin exposición.
%%% [R29] Reporte T2-T5 (sec. 1.1): umbral de exposicion 1e-8 y promedios solo sobre meses con exposicion (compute_metrics.m).
%%% [fin C29]

    \item \textbf{Supervivencia y cómputo:} número de incumplimientos de $\delta$ fuera de muestra, mínimo factor de crecimiento realizado, meses sin optimización por fallos del solucionador y fracción de meses con solución óptima.
%%% [R19b] Texto anterior: ... fallos del solucionador y fracción de fechas factibles. Se renombra porque la columna mide la fracción de meses con estado Solved, no factibilidad teórica (reporte T2-T5, C19).
\end{itemize}

\subsubsection{Inferencia sobre las diferencias}

La riqueza final es una única observación acumulada y no basta para afirmar que una estrategia es superior. Por eso la comparación estadística se basó en las diferencias mensuales pareadas del crecimiento logarítmico entre dos estrategias $A$ y $B$,
\begin{equation}\label{eq:D}
D^{A,B}_k=\log G^{A,\mathrm{real}}_{k+1}-\log G^{B,\mathrm{real}}_{k+1}.
\end{equation}
Se compararon (i) Wasserstein--Kelly con geometría~1 frente al Kelly nominal, (ii) geometría~2 frente a geometría~1 (el contraste central del estudio) y (iii) 
%%% [C30/C34] Dir. p.9 y p.11 (procedimiento de inferencia concreto; elegir entre bootstrap móvil y estacionario). REVISAR: ℓ=6 y B=10 000 son reglas a priori propuestas; se usa Künsch (1989) y se retira Politis y Romano (1994).
%%% Texto anterior: geometría~2 frente al Kelly nominal.% \textbf{[PENDIENTE: implementar el bootstrap móvil por bloques \citep{kunsch1989} --bloques de seis meses, remuestreando los pares de resultados mensuales y no cada estrategia por separado-- y reportar la media y la mediana de $D_k$, un intervalo del $95\,\%$ y la proporción de meses con $D_k>0$. Hoy solo se calculan las diferencias puntuales.]}
geometría~2 frente al Kelly nominal, cada uno con parámetros emparejados y con las configuraciones seleccionadas (Sección~\ref{sec:contrastes}); solo el contraste emparejado aísla el efecto de la geometría.

Para cuantificar la incertidumbre de la media de $D_k$ se usó un bootstrap móvil por bloques \citep{kunsch1989}: se remuestrearon bloques solapados de $\ell=6$ meses de la serie de pares $(D_k)$, preservando la alineación temporal (los resultados de las dos estrategias del mismo mes se mueven juntos), y se calcularon $B=10\,000$ réplicas de la media de $D_k$; el intervalo del $95\,\%$ es el percentil $2.5$--$97.5$ de esas medias. Se reportaron además la mediana de $D_k$ y la proporción de meses con $D_k>0$. La longitud de bloque es una regla a priori y la inferencia es sensible a ella. 
%%% [R30] Resumen T10 (T6, bootstrap_block_diff.m y test_simulado.csv): se reemplaza el PENDIENTE de implementacion y sensibilidad. Texto anterior: [PENDIENTE: implementar el procedimiento (hoy solo se calculan las diferencias puntuales) y comprobar la sensibilidad con l=3 y l=12.]
El procedimiento se implementó en MATLAB con semilla fija (20261019) y se verificó con series simuladas de media conocida y errores AR(1) de coeficiente $0.3$, con $T=84$ y 500 simulaciones, evaluando la cobertura empírica del intervalo del $95\,\%$ y el sesgo de la media. La sensibilidad se evaluó con $\ell=3$ y $\ell=12$ y se reporta junto con el caso base. Como algunos meses tienen $|D_k|<10^{-6}$ (empates numéricos), la proporción de meses con $D_k>0$ se acompaña del número de esos empates, que solo se reportan. Las comparaciones fueron los siete contrastes de la fase de prueba (Sección~\ref{sec:contrastes}).
%%% [fin R30]
%%% [fin C30/C34]

 Un intervalo que contenga el cero se interpretó como ausencia de evidencia suficiente de una diferencia media en el periodo disponible, no como prueba de equivalencia. 
%%% [C31] Dir. p.9 (a qué mes corresponde cada señal y umbral). REVISAR: yo definí la mediana sobre la ventana D_k; confirma o cambia.
%%% Texto anterior: Además, se desarrolló un desglose por régimen, meses con VIX por encima y por debajo de la mediana y meses con retorno negativo de SPY.
Como análisis descriptivo, se desglosaron las diferencias $D_{k}$ según el estado observable en la fecha de decisión $k$: meses $k+1$ cuyo VIX al cierre de $k$ estuvo por encima o por debajo de la mediana del VIX de la ventana $\mathcal D_k$ (calculada solo con datos hasta $k$, sin mirar el periodo de prueba) y meses $k+1$ precedidos por un retorno negativo de SPY en $k$. La ventana $\mathcal D_k$ comprende los 60 meses que terminan en $k$, incluido $k$; un VIX igual a la mediana se clasificó como bajo y un retorno de SPY igual a cero como no negativo. Además de cada criterio por separado se cruzaron ambos. Al usar información disponible en $k$, el desglose no depende del resultado realizado del mes evaluado; aun así es descriptivo y no prueba capacidad de anticipación.
%%% [R31] Resumen T10 (T7, results_regimen/): confirmado que el desglose se calculo como dice el texto (informacion disponible en k, mediana de la ventana que incluye k). Se agregan empates, cruce y tamanos de grupo. Texto anterior: [PENDIENTE: confirmar que este desglose se calculó así.]
%%% [fin C31]

 % \textbf{[PENDIENTE (opcional): desglose por régimen --meses con VIX por encima y por debajo de la mediana y meses con retorno negativo de SPY--, solo descriptivo.]}

\subsection{Implementación y verificación}\label{sec:impl}

Los datos se descargaron y se transformaron a frecuencia mensual con Python (\texttt{yfinance}); el modelo se implementó en MATLAB (R2026a) con CVX 2.2 \citep{grantboyd} y el solucionador MOSEK 9.1.9 \citep{mosek}, que admite de forma nativa las restricciones del cono exponencial presentes en \eqref{eq:nominal} y \eqref{eq:dual}. 
%%% [C32] Dir. p.9 (equivalencia numérica con η=0: tolerancia explícita y soluciones múltiples). REVISAR: coloca las tolerancias de test_wdro_nominal_equivalence.m.
%%% Texto anterior: %%% [R32] Reporte T2-T5 (sec. 4 y 5/C32): tolerancias y resultados de las verificaciones (a), (b) y (c). Texto anterior (comentado, incluye los PENDIENTE de tolerancias y repositorio):
%%%   Antes de interpretar resultados se verificó que (a) con $\rho_k=0$ el modelo Wasserstein--Kelly reproduce la solución del modelo nominal, tanto en una instancia pequeña como en el backtest completo;
%%%   Antes de interpretar resultados se verificó que (a) con $\rho_k=0$ el modelo Wasserstein--Kelly reproduce la solución del modelo nominal, tanto en una instancia pequeña como en el backtest completo (como pueden existir soluciones óptimas múltiples, se compararon los valores objetivos y la factibilidad con una tolerancia explícita, 
%%%   %\textbf{[PENDIENTE: valor de la tolerancia]}
%%%   ,y las ponderaciones solo cuando la solución es única o estable, con otra tolerancia %\textbf{[PENDIENTE: valor]}
%%%   ; un desacuerdo menor entre variables con igual valor objetivo no se interpretó como error del modelo);
%%%   %%% [fin C32]
%%%    (b) las posiciones cumplen el balance $x_{0k}+\sum_j x_{jk}=1-\tau\,\mathrm{TO}_k$, los límites por activo y la cota de supervivencia sobre la ventana; y (c) con $\eta=0$ las dos geometrías coinciden.
%%%   %\textbf{[PENDIENTE: confirmar que cada una de estas verificaciones se corrió y quedó registrada; retirar la que no.]} \textbf{[PENDIENTE: indicar el repositorio con el código y los datos originales para reproducibilidad.]}
Antes de interpretar resultados se verificó, con tolerancias explícitas, que (a) con $\rho_k=0$ el modelo Wasserstein--Kelly reproduce la solución del modelo nominal, tanto en tres instancias individuales (fechas de decisión de diciembre de 2014, diciembre de 2018 y junio de 2022, partiendo de efectivo y con $\delta=0.75$) como en el backtest completo de 132 meses para los seis valores de $\delta$. Como pueden existir soluciones óptimas múltiples, se compararon los valores objetivos con una tolerancia absoluta de $10^{-6}$, la factibilidad con un residuo máximo de $10^{-7}$ en las restricciones y los factores de crecimiento mensuales del backtest con tolerancia $10^{-6}$; las ponderaciones $(x_0,x)$ se compararon solo cuando la solución es única, con tolerancia $10^{-5}$. La unicidad se evaluó así: el objetivo es estrictamente cóncavo en el vector de factores de crecimiento $G=A\,(x_0,x)$, donde $A$ tiene por columnas los factores de crecimiento del activo de bajo riesgo y de las acciones en cada escenario, por lo que $G$ es único y $(x_0,x)$ también lo es si $A$ tiene rango completo por columnas; en las tres instancias se calculó el rango y el número de condición de $A$. Un desacuerdo menor entre variables con igual valor objetivo no se habría interpretado como error del modelo. (b) Las posiciones cumplen el balance $x_{0k}+\sum_j x_{jk}=1-\tau\,\mathrm{TO}_k$, los límites por activo y la cota de supervivencia sobre la ventana, comprobado en todos los meses guardados de las 66 celdas de la geometría~1 (malla completa) y las 66 de la geometría~2 (validación). (c) Con $\eta=0$ las dos geometrías coinciden: se compararon sus trayectorias en los 48 meses de validación y para los seis valores de $\delta$, con tolerancia $10^{-8}$.
%%% [fin R32]
%\textbf{[PENDIENTE: indicar el repositorio con el código y los datos originales para reproducibilidad.]}

\subsection{Limitaciones del diseño}\label{sec:limit}

\begin{itemize}
    \item El universo es una lista fija de empresas grandes que sobrevivieron hasta hoy, por lo que no elimina el sesgo de supervivencia. Los resultados son una comparación metodológica sobre ese universo, no una reconstrucción del conjunto de oportunidades de 2010.
    \item El backtest es una única trayectoria histórica; la validación y la prueba separan la elección de parámetros de la evaluación, pero no sustituyen la evidencia de varios universos o periodos.
    \item Los costos proporcionales de $0.001$ no reflejan comisiones ni diferenciales reales de cada activo.
    \item El rendimiento de DGS3MO es un proxy del retorno de bajo riesgo.
    \item El soporte de la bola es finito: el adversario solo redistribuye probabilidad entre meses ya observados y no puede crear escenarios nuevos.
    \item Los pesos $\omega$ de la geometría~2 se fijaron y no se calibraron; el VIX entra solo por su nivel.
    
%%% [C10/C15/C23b] Dir. p.5, p.7 y p.8: se agregan a las limitaciones la equiprobabilidad, el riesgo de sobreajuste de la validación y la dependencia de las reglas a priori.
\item La cota $\delta$ protege únicamente frente a los escenarios de la ventana.
    \item La ventana móvil trata los $N$ meses como equiprobables; esto es la distribución empírica de referencia y no supone que los retornos sean independientes, estacionarios ni idénticamente distribuidos, por lo que un cambio de estado dentro de la ventana no se refleja en los pesos.
    \item La selección se basa en la riqueza final sobre 48 meses de validación y 60 combinaciones elegibles $(\delta,\eta)$ por geometría (de 66 exploradas), lo que puede sobreajustar la validación; la configuración elegida no es la mejor en sentido multicriterio.
    \item Los pesos de la geometría~2 y la longitud de bloque del bootstrap son reglas a priori, y las conclusiones están condicionadas a ellas.
%%% [C7c] Auditoría: se agrega la limitación de fuente única y del cierre de diciembre de 2025.
    \item Los precios provienen de una única fuente oficial (Yahoo Finance), contrastada contra Tiingo y CBOE con 4 discrepancias sin explicar; además, el cierre de diciembre de 2025 corresponde al 30 de diciembre y no al 31.
%%% [R-T9] Reportes T9 y de prueba: limitaciones de proceso de la selección y del uso previo del periodo de prueba.
    \item La regla de selección se fijó después de observar la tabla de validación (se excluyó $\eta=0$ por esa razón), y su tolerancia de empate se calibró con el ruido numérico medido en la propia validación.
%%% [R-T9b] Aclaración del autor: las corridas previas fueron de verificación de la optimización y de exploración de factibilidad; la malla se construyó con ese criterio. Texto anterior: (lista de resultados de prueba ya calculados, con REVISAR)
    \item Antes de fijar la regla de selección se habían ejecutado, sobre la trayectoria completa (que incluye el periodo de prueba), corridas cuyo propósito era comprobar que la optimización funcionaba y explorar en qué regiones de $(\delta,\eta)$ había factibilidad; la malla de la Tabla~\ref{tab:param} se construyó con ese criterio y no con el desempeño. Ningún parámetro se eligió con los resultados de prueba, pues la regla de selección usa solo validación; aun así, esas corridas produjeron métricas del periodo de prueba, por lo que este no es estrictamente virgen.
\end{itemize}
%%% [fin C10/C15/C23b]

%%% [C33] Dir. p.10 (bibliografía): se unifican las etiquetas autor-año en español («y» en lugar de «and») y las iniciales de Li. Pendiente de revisión manual: formato de las referencias técnicas de CVX/MOSEK (versión y fecha de consulta).
%%% [R33] Reporte T2-T5: se agregan las versiones (CVX 2.2 build 1148; MOSEK 9.1.9). Sigue pendiente la fecha de consulta de ambas referencias.
\begin{thebibliography}{99}


\bibitem[Markowitz(1952)]{markowitz1952}
H.~Markowitz.
\newblock Portfolio Selection.
\newblock \emph{The Journal of Finance}, 7(1):77--91, 1952.
\newblock DOI: 10.1111/j.1540-6261.1952.tb01525.x.

\bibitem[Rockafellar y Uryasev(2000)]{rockafellaruryasev2000}
R.~T. Rockafellar y S.~Uryasev.
\newblock Optimization of Conditional Value-at-Risk.
\newblock \emph{Journal of Risk}, 2(3):21--41, 2000.

\bibitem[Kelly(1956)]{kelly1956}
J.~L. Kelly Jr.
\newblock A New Interpretation of Information Rate.
\newblock \emph{Bell System Technical Journal}, 35(4):917--926, 1956.
\newblock DOI: 10.1002/j.1538-7305.1956.tb03809.x.

\bibitem[Breiman(1961)]{breiman1961}
L.~Breiman.
\newblock Optimal Gambling Systems for Favorable Games.
\newblock En \emph{Proceedings of the Fourth Berkeley Symposium on
  Mathematical Statistics and Probability}, vol.~1, pp.~65--78, 1961.

\bibitem[Sun y Boyd(2018)]{sunboyd2018}
Q.~Sun y S.~Boyd.
\newblock Distributional Robust Kelly Gambling: Optimal Strategy under
  Uncertainty in the Long Run.
\newblock arXiv:1812.10371, 2018.
\newblock \url{https://arxiv.org/abs/1812.10371}.

\bibitem[Whaley(2000)]{whaley2000}
R.~E. Whaley.
\newblock The Investor Fear Gauge.
\newblock \emph{The Journal of Portfolio Management}, 26(3):12--17, 2000.
\newblock DOI: 10.3905/jpm.2000.319728.

\bibitem[French et al.(1987)]{french1987}
K.~R. French, G.~W. Schwert y R.~F. Stambaugh.
\newblock Expected Stock Returns and Volatility.
\newblock \emph{Journal of Financial Economics}, 19(1):3--29, 1987.
\newblock DOI: 10.1016/0304-405X(87)90026-2.

\bibitem[Li(2023)]{li2023}
J.~Y.-M. Li.
\newblock Wasserstein--Kelly Portfolios: A Robust Data-Driven Solution to
  Optimize Portfolio Growth.
\newblock arXiv:2302.13979, 2023.
\newblock \url{https://arxiv.org/abs/2302.13979}.

\bibitem[Hsieh y Yu(2024)]{hsiehyu2024}
C.-H. Hsieh y X.-R. Yu.
\newblock On Cost-Sensitive Distributionally Robust Log-Optimal Portfolio.
\newblock arXiv:2410.23536, 2024.
\newblock \url{https://arxiv.org/abs/2410.23536}.

\bibitem[Liu et al.(2022)]{liu2022}
J.~Liu, J.~Wu, B.~Li y P.~Cui.
\newblock Distributionally Robust Optimization with Data Geometry.
\newblock En \emph{Advances in Neural Information Processing Systems 35},
2022.
\newblock \url{https://proceedings.neurips.cc/paper_files/paper/2022/hash/da535999561b932f56efdd559498282e-Abstract-Conference.html}.

\bibitem[Mohajerin Esfahani y Kuhn(2018)]{mohajerin2018}
 P.~Mohajerin Esfahani y D.~Kuhn.
\newblock Data-driven distributionally robust optimization using the
   Wasserstein metric: performance guarantees and tractable reformulations.
\newblock \emph{Mathematical Programming}, 171:115--166, 2018.
\newblock DOI: 10.1007/s10107-017-1172-1.

 \bibitem[Kuhn et al.(2019)]{kuhn2019}
 D.~Kuhn, P.~Mohajerin Esfahani, V.~A. Nguyen y S.~Shafieezadeh-Abadeh.
 \newblock Wasserstein Distributionally Robust Optimization: Theory and
  Applications in Machine Learning.
\newblock En \emph{INFORMS TutORials in Operations Research}, pp.~130--166, 2019.


%%% [C2b] Faltaba la barra invertida en \bibitem (la entrada no se compilaba).
\bibitem[Sun y Zou(2025)]{sunzou2025}
Y.~Sun y Z.~Zou.
\newblock Data-driven distributionally robust Kelly portfolio optimization
  based on coherent Wasserstein metrics.
\newblock \emph{Journal of University of Science and Technology of China
  (JUSTC)}, 55(8):0805-1--0805-11, 2025.
\newblock DOI: 10.52396/JUSTC-2024-0030.

 \bibitem[Grant y Boyd(s.f.)]{grantboyd}
 M.~Grant y S.~Boyd.
 \newblock Using MOSEK with CVX.
 \newblock CVX Users' Guide, versión 2.2 (build 1148). 
 \url{https://cvxr.com/cvx/doc/mosek.html}.
\newblock Consultado el 15 de octubre de 2026.

\bibitem[Sharpe(1994)]{sharpe1994}
W.~F. Sharpe.
\newblock The Sharpe Ratio.
\newblock \emph{The Journal of Portfolio Management}, 21(1):49--58, 1994.
\newblock DOI: 10.3905/jpm.1994.409501.

\bibitem[Sortino y Price(1994)]{sortinoprice1994}
F.~A. Sortino y L.~N. Price.
\newblock Performance Measurement in a Downside Risk Framework.
\newblock \emph{The Journal of Investing}, 3(3):59--64, 1994.
\newblock DOI: 10.3905/joi.3.3.59.

\bibitem[Young(1991)]{young1991}
T.~W. Young.
\newblock Calmar Ratio: A Smoother Tool.
\newblock \emph{Futures}, 20(1):40, 1991.

\bibitem[K{\"u}nsch(1989)]{kunsch1989}
H.~R. K{\"u}nsch.
\newblock The Jackknife and the Bootstrap for General Stationary Observations.
\newblock \emph{The Annals of Statistics}, 17(3):1217--1241, 1989.
\newblock DOI: 10.1214/aos/1176347265.

%%% [C34] Dir. p.11: se decidió el bootstrap MÓVIL por bloques (Künsch, 1989); Politis y Romano (1994) corresponde al bootstrap estacionario y no se usa, por lo que se retira de la bibliografía.
%% \bibitem[Politis y Romano(1994)]{politisromano1994}
%% D.~N. Politis y J.~P. Romano.
%% \newblock The Stationary Bootstrap.
%% \newblock \emph{Journal of the American Statistical Association}, 89(428):1303--1313, 1994.
%% \newblock DOI: 10.1080/01621459.1994.10476870.

\bibitem[MOSEK ApS(s.f.)]{mosek}
MOSEK ApS.
\newblock \emph{MOSEK Optimization Toolbox for MATLAB}, versión 9.1.9.
\newblock \url{https://docs.mosek.com}.
\newblock Consultado el 15 de octubre de 2026.
\end{thebibliography}

\end{document}